% ==============================================================================
% 250_FFGFT_Information_SL_En_ch.tex
% Chapter content – included via wrapper
% Author: Johann Pascher · May 2026
% Revised after DeepSeek review: all 6 improvements integrated
% ==============================================================================

\section*{Abstract}

This document clarifies what \emph{information} means in FFGFT in the context of black holes. It shows that the Hawking information paradox rests on a category confusion: the classical formulation treats information as an epistemic concept. FFGFT distinguishes the \textbf{fundamental information}: the topological structure of winding numbers on $T^4$ --- comparable to a topological charge, not to a quantity of bits. This topological charge is \textbf{topologically preserved}, independently of any observer. The compactification of the time dimension on $T^4$ undermines one of the three presuppositions of the paradox, and thermodynamic irreversibility is fully compatible with topological information preservation.

\tableofcontents
\newpage

% ------------------------------------------------------------------------------
\section{The Hawking Information Paradox and Its Presuppositions}

Stephen Hawking showed in 1974 that black holes emit thermal radiation \cite{Hawking1974}. This radiation carries no information about the specific state of the infalling matter. If the black hole evaporates completely, the original quantum information would be destroyed --- violating unitarity.

The paradox presupposes three conditions that do not all hold in FFGFT:

\begin{enumerate}
  \item Information is the complete description of the quantum state of all infalling particles.
  \item Time is continuous and unidirectionally oriented --- what fell into the black hole ``in the past'' is ``forever'' inaccessible.
  \item The classical singularity is physically real --- the state of matter is destroyed there.
\end{enumerate}

\begin{wichtig}[Why does the paradox arise at all?]
Hawking (implicitly) presupposes that information is \emph{epistemic} --- what an observer can reconstruct --- and that time is continuous. Neither is fundamental in FFGFT: information is a topological charge, not an epistemic quantity; and time is compactified on $T^4$, not continuous. The paradox dissolves not by adding new physics, but by clarifying the presuppositions.
\end{wichtig}

FFGFT rejects all three presuppositions in specific ways --- not by assertion, but through the geometric structure of the framework.

% ------------------------------------------------------------------------------
\section{What Information in FFGFT Is \emph{Not}}

\subsection{Not our description of the black hole}

Information in the FFGFT sense is \textbf{not} our knowledge or understanding of the black hole. The Hawking paradox is an \textbf{epistemic} problem: it asks whether an observer outside the horizon can fully reconstruct the original quantum state. FFGFT answers a different question: does the fundamental topological structure continue to exist?

\begin{wichtig}
Information in the sense of FFGFT is not an epistemic category (what we can know), but an \textbf{ontological} category (what structurally exists). The inability of an external observer to reconstruct information is an access problem --- not an existence problem.

\textbf{Observable consequence:} If the topological information were truly destroyed, the Hawking radiation would be exactly black-body radiation. The modified spectrum (due to the $T^4$ winding structure) is the falsifiable difference --- this empirically grounds the access problem.
\end{wichtig}

\subsection{Not the sum of all detail states}

Information in the FFGFT sense is also not the sum of all information required to fully describe all particles. These detail states can be epistemically lost through thermalisation. That is thermodynamic irreversibility, not ontological information destruction.

% ------------------------------------------------------------------------------
\section{What Information in FFGFT \emph{Is}: Fundamental Information}

\subsection{Topological invariants as fundamental information}

In FFGFT, particles are winding configurations on the 4D identification torus $T^4$ with quantum numbers $(n_\theta, n_\varphi, n_3, n_4)$. These winding numbers are \textbf{topological invariants}.

\begin{keyresult}[Definition of Fundamental Information]
The fundamental information of a physical system in FFGFT is the \textbf{topological charge} of its constituents: the winding numbers $(n_\theta, n_\varphi, n_3, n_4)$ on $T^4$. These define \emph{what kind} of particles exist --- not \emph{where} they are, not \emph{how} they behave, not \emph{in what exact quantum state} they find themselves.

Fundamental information is not a \emph{quantity of bits} (Shannon), but a \textbf{topological charge} --- comparable to a magnetic monopole charge or a winding number in string theory. This immediately decouples FFGFT from the usual bit-oriented information paradox debate.
\end{keyresult}

Topological invariants cannot be reduced to zero through continuous transformations. A state with winding number $n = 1$ cannot be continuously deformed into $n = 0$ without a topological discontinuity --- with measurable physical consequences.

\begin{erkenntnis}
\textbf{No destruction operator acts on topological invariants.} This is the central statement: there is no physical process that can reduce winding numbers to zero without generating a measurable discontinuity. Classical singularities would be such a process --- which is why they are structurally excluded in FFGFT.
\end{erkenntnis}

\subsection{The foundational relation as structure preservation}

The fundamental relation $\tilde{T} \cdot m = 1$ connects the T-field and mass at every point. In the limit of extreme compression:
\begin{equation}
  \tilde{T} \rightarrow L_0/c = \xipar \cdot t_P > 0
\end{equation}
The T-field never reaches zero --- the minimum value is $L_0/c$. The winding structure of the field configuration is preserved.

\subsection{Algebraic protection against singularities (Dok.~078)}

Here lies the most elegant proof --- not by a postulate, but by the algebraic structure of the foundational relation itself.

From $\tilde{T} \cdot m = 1$:
\begin{equation}
  \tilde{T} = \frac{1}{m}
\end{equation}

If $m \rightarrow \infty$, then $\tilde{T} \rightarrow 0$. But reaching $\tilde{T} = 0$ exactly would require $m = \infty$ exactly --- algebraically impossible.

\begin{keyresult}[Algebraic Protection --- Dok.~078]
$\tilde{T} \cdot m = 1$ acts as algebraic protection against singularities:
\begin{itemize}
  \item $m \rightarrow \infty$ implies $\tilde{T} \rightarrow 0$ --- but only \textbf{asymptotically}, never exactly.
  \item $\tilde{T} = 0$ exactly would require $m = \infty$ exactly --- algebraically impossible.
  \item The singularity is not a forbidden state, but an \textbf{unreachable limit}.
\end{itemize}
This is not a geometric and not a quantum mechanical argument --- it is an \textbf{algebraic argument} from the foundational relation alone. It holds independently of the size of the black hole and regardless of how far one drives the compression.
\end{keyresult}

\subsection{The interior operator: what happens at $L_0$}

What happens at the boundary $L_0$ in the interior of a black hole is not fully open in FFGFT.

\textbf{The T-field at its minimum.}
$\tilde{T}_\text{min} = \xipar \cdot t_P > 0$ --- a finite, positive, non-singular value.

\textbf{Maximum winding density.}
One winding quantum per $L_0^4$ spacetime volume --- analogous to the lower bound in a lattice with finite volume. This is the densest stable configuration on the $\xipar$ ladder (level $N = 0$). Below this level the state space ends.

\textbf{Lagrangian stiffness.}
The term $-\frac{1}{2}m_T^2(\partial m)^2$ with $m_T = \xipar/\ell$ (Dok.~201) prevents $\partial m$ from diverging --- algebraic property of the field equations, not an external force.

\textbf{Topological bounce condition.}
No stable configurations below $N = 0$ $\Rightarrow$ no states the system can transition into. This is a \textbf{state space boundary}, not a force-counterforce mechanism.

\textbf{Time in the interior.}
$\tilde{T} \rightarrow L_0/c$: local time quantum at its minimum. Time does not ``freeze'' ($\tilde{T} = 0$ would be required), but the local time structure reaches its smallest possible value. The modified Hawking spectrum arises precisely at this boundary.

\textbf{Formal criterion for the transition operator.}
An interior operator would need to describe the mapping of an incoming winding configuration $(n_i)$ onto the extremal configuration at $L_0$ --- under preservation of topological charge, with maximum winding density. The $T^4$ algebra suggests this operator is \textbf{unitary}, but commutes non-trivially with the thermodynamic limit. That is the open research question --- precisely stated, not a fundamental gap.

\begin{erkenntnis}[Interior of a Black Hole in FFGFT]
The interior is not unknown territory, but an extreme state:
\begin{itemize}
  \item $\tilde{T}_\text{min} = \xipar \cdot t_P > 0$ --- finite, not singular
  \item Maximum winding density: $1/L_0^4$ per spacetime volume
  \item Lagrangian stiffness: $\partial m$ does not diverge
  \item State space boundary at $N = 0$
  \item Unitary transition operator (details open)
\end{itemize}
\end{erkenntnis}

\subsection{Fundamental information is not the description of the totality}

The fundamental information is the \textbf{class} of configurations (topological charge), not their \textbf{instance} (quantum state):
\begin{itemize}
  \item Class: ``There is an electron'' (winding numbers defined)
  \item Instance: ``The electron is at position $x$ with spin $+\frac{1}{2}$'' (quantum state)
\end{itemize}
FFGFT preserves the class. The instance can become thermodynamically irreversible.

% ------------------------------------------------------------------------------
\section{Black Holes as Inversion --- Not as Destruction}

\subsection{The topological structure remains, inverted}

A black hole is not a destruction state, but an \textbf{inversion} of normal matter structure: winding configurations maximally compressed, T-field at minimum $L_0/c$. What fell in as an electron remains topologically classifiable as an electron --- in an extreme field regime.

\subsection{Irreversibility is thermodynamic, not ontological}

\begin{erkenntnis}
Irreversibility and information preservation operate on different levels:
\begin{itemize}
  \item \textbf{Thermodynamic:} Directed, irreversible, entropy increases.
  \item \textbf{Topological:} Winding numbers \textbf{topologically preserved}. No destruction operator acts on topological invariants.
\end{itemize}
The Hawking paradox arises from confusing these two levels.
\end{erkenntnis}

% ------------------------------------------------------------------------------
\section{The Compactified Time Dimension}

\subsection{Time is not continuous --- it is compactified}

In FFGFT, the time dimension on $T^4$ is \textbf{compactified} --- the winding number $n_4$ is a topological invariant like the others.

\begin{wichtig}
Compactification means \textbf{periodicity of field configurations} in the time direction --- not reversibility of the thermodynamic time. A macroscopic process remains irreversible. The $T^4$ topology is time-direction-neutral; the thermodynamic time direction is emergent, not fundamental.
\end{wichtig}

\subsection{What compactification changes}

The statement ``information is forever lost in the past'' presupposes an absolute ontological past boundary. In a compactified time dimension, this absoluteness is not self-evident: the topological field structure is not referenced to the observer's time direction. $n_4$ remains topologically preserved.

\subsection{Our experienced time is unambiguously directed}

The experienced time direction is a thermodynamic phenomenon (entropy increase), not a geometric primitive. In FFGFT, the time direction is emergent, not fundamental --- compactification operates at the topological level and does not mystically overload the theory.

% ------------------------------------------------------------------------------
\section{What FFGFT Concretely Claims --- and What It Does Not}

\begin{enumerate}
  \item \textbf{No classical singularity:} $L_0 = \xipar \cdot \ell_P$ as absolute lower density bound.
  \item \textbf{Topological fundamental information (charge) preserved:} Winding numbers as topological invariants, algebraically protected by $\tilde{T} \cdot m = 1$ (Dok.~078).
  \item \textbf{Modified Hawking spectrum:} Falsifiable prediction --- deviation from black-body spectrum due to $T^4$ winding structure.
  \item \textbf{Gravitational wave echoes:} After mergers as a further test.
\end{enumerate}

\textbf{Not claimed:} Complete state reconstruction; carrier recovery; fully specified transition operator near $N = 0$; suspension of thermodynamic irreversibility.

\begin{keyresult}[Precise FFGFT Statement]
FFGFT claims: the \textbf{topological fundamental charge} --- the winding numbers of the infalling matter --- is \textbf{topologically preserved}. It is not reconstructible by an external observer (epistemic access problem), but it continues to exist as a geometric property of the $T^4$ field.

If it were truly destroyed, the Hawking spectrum would be exactly thermal --- the modified spectrum is the measurable difference.
\end{keyresult}

% ------------------------------------------------------------------------------
\section{Connection to the $\xipar$ Recursion}

The $\xipar$ ladder defines which winding configurations on $T^4$ are stable. Black holes represent configurations at the lower end ($N = 0$) --- within the $\xipar$ structure, not outside it. The ``inverted'' configuration of a black hole is also part of the $\xipar$ ladder and topologically well-determined.

% ------------------------------------------------------------------------------
\section*{Summary}

\begin{longtable}{>{\bfseries}p{6cm} p{8cm}}
\toprule
Question & FFGFT Answer \\
\midrule
Why does the paradox arise? & Hawking presupposes information as epistemic and time as continuous --- neither fundamental in FFGFT \\
What is fundamental information? & Topological charge: winding numbers $(n_\theta, n_\varphi, n_3, n_4)$ on $T^4$ --- not bits \\
What is not information? & Our knowledge of the state; carrier identity; complete state description \\
Is information lost? & Topological charge: no (topologically preserved). Detail state: access problem \\
Empirical grounding & Modified Hawking spectrum: if information were destroyed, radiation would be exactly thermal \\
Algebraic protection (Dok.~078) & $m \rightarrow \infty$ implies $\tilde{T} \rightarrow 0$ only asymptotically --- singularity algebraically unreachable \\
What happens at $L_0$? & $\tilde{T}_\text{min} > 0$ · Max.\ winding density · Lagrangian stiffness · State space boundary $N=0$ \\
Is there a singularity? & No --- $L_0$ as lower bound; state space ends at $N = 0$ \\
Is the process irreversible? & Thermodynamically yes. Topologically: fundamental structure remains \\
Time compactification & Periodicity of field configurations, not reversal of thermodynamic time \\
Falsifiable tests & Hawking spectrum (not exactly thermal); gravitational wave echoes \\
What remains open? & Transition operator $(n_i) \rightarrow$ extremal configuration at $L_0$ --- unitary, details open \\
\bottomrule
\end{longtable}
